% ============================================================
%  出版说明
% ============================================================
\section*{出版说明}
\addcontentsline{toc}{section}{出版说明}
\markboth{出版说明}{出版说明}

这本小书是一次重排。它的底本不是论文，而是 2021 年 9 月 23 日 Edward Witten 在剑桥 Isaac Newton Institute 所作的一场演讲的 54 张幻灯片。幻灯片每页只有一两句话，公式与示意图各自独立成页；照着读固然可以，却不像一本书。我们按演讲本身的叙事顺序，把这些要点连缀成连续的文字，译成中文，并在关键处补上技术性的批注与史实考订，使它既是一份可靠的译本，也可以当作进入“电磁对偶—几何 Langlands”这条线索的导论。

\medskip
\noindent\textbf{一、底本。}\quad
底本文件为 \emph{Duality, the Langlands Correspondence, and the Influence of Michael Atiyah}，作者 Edward Witten，54 页，页面尺寸 362.835\,pt $\times$ 272.126\,pt（16:9），由 \LaTeX{} 的 Beamer 文档类生成；首页注明演讲场合为 Isaac Newton Institute, September 23, 2021。文件大小 2\,006\,439 字节，SHA-256 校验值见“工程文件/底本说明.md”。底本 PDF 因版权与体积原因不收入仓库，正文与批注中的页码一律指这 54 张幻灯片。

\medskip
\noindent\textbf{二、译本的分工。}\quad
本书的文字分四层，版式上彼此可辨：

\begin{itemize}
  \item \textbf{正文}（黑色）：Witten 演讲的忠实翻译。原句中的强调、犹豫与口语——\emph{wild}、\emph{I doubt}、\emph{smoking gun}——尽量保留语气，不改成四平八稳的书面语。
  \item \textbf{译注}（橙色框）：术语的选择、译法的理由、原文歧义处的处理。编号按章，如“译注 3.2”。
  \item \textbf{批注}（蓝色框）：技术性补充与史实考订，为的是让不熟悉这一行的读者也能跟上。批注是译者所加，不属于原文；只想读 Witten 的读者可以整段跳过。编号按章，如“批注 5.1”。
  \item \textbf{文献摘录}（灰色框）：底本中有五页直接贴了论文首页或摘要的扫描图（GNO、Montonen–Olive、Witten–Olive、Osborn、Sen）。我们不搬扫描图，而是把摘要原文排印出来，附中译。
\end{itemize}

正文中形如 \srcp{12} 的灰色小角标，标出这一段话在底本中起始于第几页幻灯片，便于对照回查。全书的页—段对应表见“工程文件/底本对应表.tsv”。

\medskip
\noindent\textbf{三、译名。}\quad
译名依数译术语库：\emph{Langlands 对偶群}、\emph{几何 Langlands 对应}、\emph{Hecke 特征层}、\emph{几何 Satake 等价}、\emph{支配余权}、\emph{余根}、\emph{单极子}、\emph{瞬子}、\emph{Bogomolny 方程}。术语首次出现时给出英文，附录 A 有完整对照表。人名与机构名保留拉丁字母，不音译，这是数译的惯例。

\medskip
\noindent\textbf{四、关于原文。}\quad
我们尽可能保留原文的省略与保留：Witten 说“粗略地说”“我相信”“我不确定”的地方，译文不作补全；凡原文未写、而由我们补出的数学内容，一律放进批注，并在批注里说明这是补出的。底本中确有前后不一致之处（例如第 3 章耦合常数的写法，与今天通行的 $\tau \to -1/\tau$ 差一个因子），我们照原文译出，另在批注中交代。讲稿的口吻属于 2021 年；2021 年之后数学与物理的进展，只在批注中以“补记”标出，不改动正文的时代感。

\medskip
\noindent\textbf{五、责任。}\quad
原文版权归 Edward Witten 所有。本译本与批注由数译以 CC BY-NC-SA 4.0 发布，仅供学习与研究；引用原文之处均为学术评论。译文错误由译者负责；批注中的公式与史实，凡有可疑，均在文中标明，欢迎指正。
